\label{anlage:transkript-bunschoten}

\begin{Table}[title=Interviewinformationen]
\begin{SemioTable}{@{}T{0.3\linewidth-2\tabcolsep}T{0.7\linewidth-2\tabcolsep}@{}}
\SemioTableRow{Interviewpartner & Raoul Bunschoten}
\SemioTableRow{Interviewer & Kinan Sarakbi}
\SemioTableRow{Datum & 04.06.2026}
\SemioTableRow{Ort / Format & Einstein Center Digital Future}
\SemioTableRow{Sprache & Englisch}
\SemioTableRow{Transkripttyp & Bearbeitetes Transkript}
\end{SemioTable}
\end{Table}

Dieses Transkript wurde zur besseren Lesbarkeit sprachlich überarbeitet. Füllwörter, Wortwiederholungen, kurze Bestätigungen, Husten sowie Hintergrundgeräusche wurden entfernt. Satzstrukturen und offensichtliche Fehler der automatischen Transkription wurden dort korrigiert, wo der beabsichtigte Sinn eindeutig war. Fachliche Inhalte, Beispiele und Begriffe wurden beibehalten. Ein kurzer persönlicher Gesprächsteil wurde ausgelassen und entsprechend gekennzeichnet. Da es sich um ein bearbeitetes und nicht um ein wortgetreues Transkript handelt, sollten direkte Zitate nach Möglichkeit mit der Originalaufnahme abgeglichen werden.

\subsection{Fragmentierte Wertschöpfungsketten und digitale Kontinuität}

\textbf{Interviewer:} How do you model a workflow that has traditionally been disconnected, and how does the VCDM bring fragmented sectors into a connected digital value chain?

\textbf{Raoul Bunschoten:} The original premise of Bauhütte 4.0 was that there is a workflow of bio-technic urban production. That was one of the premises of the programme we began seven years ago. However, this workflow has problems. There are many gaps in its continuity and many pain points where activities are not connected or where different sectors work according to completely different standards.

This is where the problems begin. Suppose you want to build a timber project, such as a project at Tegel. Where does the wood come from? We are involved with this question at a practical level through Tegel Projekt. How do you know that the wood is actually what the label says it is? How do you know that it is the correct type of wood, where it came from, and whether it was sourced ethically? A great deal of the wood sold through Austria, for example, actually comes from Romania.

There is no continuity between the different sectors of what we call the value chain. At one point, I began referring to it as “forest-to-city, city-to-forest.” That has also been published. We are trying to develop a more comprehensive term for the space in which these material flows take place.

Today, requirements such as life-cycle assessment calculations, ESG criteria, building regulations, and procurement processes mean that you must be able to demonstrate how much wood is in a building, what its quality is, whether it can be reused later, and where it is located. These are questions of circularity and the future reuse of components.

We recently had Concular and Madaster involved in one of our events. When we formulated the core principles of Bio-Tech 4.0, we explicitly stated that we would address discontinuity, particularly digital discontinuity. This became a central focus. There are two broad issues: renewable materials and sustainable urban construction, and the digital tools required to know what is located where and how it can be managed.

That was part of the original concept. We approached it partly from the perspective of urbanism, based on the work I was developing through my chair, and partly from the production perspective represented by Fraunhofer. The question is: how do you create a value chain, and how do you produce a city? We therefore decided to focus on the gaps and fragments.

Over the years, this work developed into Value Chain Digital Modeling. We completed phase one and are now in phase two. The VCDM is intended to model the dynamics of the overall value chain in as many of its aspects as possible. We are not fully there yet, because the system is extremely complex. The aim is to simulate situations in which a gap might be overcome or a connection created between previously disconnected elements.

This leads directly to the issue of datasets. How do you compare datasets that are fundamentally different in nature? Many types of data exist along the value chain, and the question is how they can be combined. We therefore began by analysing standards: the standards used for different datasets, the standards applied to processes, how a machine operates, and how instructions are communicated to that machine. Architects and design are part of this, but design is only one sector within the entire value chain.

\subsection{Akteure, Kommunikation und Aushandlung}

\textbf{Interviewer:} That leads to my next question about stakeholders. Who are the principal actors involved across the value chain? When stakeholders have different data requirements or conflicting interests, how do you represent and manage those differences?

\textbf{Raoul Bunschoten:} We began by trying to identify an overall set of stakeholders covering the full value chain. That is difficult because the value chain itself must first be defined. Increasingly, I think not simply in terms of construction but in terms of urban construction. This continues the work of my chair in urbanism: how cities are made and how construction can move from individual projects toward a macro-urban scale.

The stakeholders include forest owners and forest specialists. We speak with organisations such as HOWOGE, Berliner Forsten, technical institutes, and individual forest owners. Berlin itself is a forest owner, and I recently met with people responsible for this area.

The value chain then includes producers: the people who cut and process the wood and prefabricate components. This relates to what I previously called intelligent prefabrication. After production come planning and design, including component systems and building kits. The question is how to define elements that can be produced and how production can be scaled.

The stakeholders also include inhabitants, associations, state-owned housing companies, and those responsible for building management. Management covers energy, water, and everything that passes through and supports a building during use. After that come adaptation, reuse, recycling, and the later stages of the material cycle. The result is a very broad range of stakeholders across the value chain.

The situation becomes especially interesting because these stakeholders do not only use different standards; they also speak in different ways. We have conducted workshops that address the non-technological dimensions of the process: the psychological, social, and cybernetic aspects of how people communicate. This has been part of the work of my chair for many years and is central to my book \textit{Urban Curation}. The question is how people can be brought into negotiation with one another.

Connecting the interests of different stakeholders or stakeholder groups may be a completely non-technological process. It may happen around a table, as it is happening now. Simulation and modeling tools are useful because they provide information about what exists where. They are sometimes described as digital twins, although that is a simplified description of what we are actually doing.

You need to know what is located where. When you do not know, you must also be able to speculate and test assumptions. You then simulate processes that might overcome the gaps and pain points we discussed earlier. In this sense, the VCDM is a planning-support tool. It is not intended to produce a solution. We are not interested in a design solution generated by the tool; we are interested in simulations that stakeholders can examine and follow. This is where scenario techniques become important.

\textit{Das Interview wird kurz durch einen Telefonanruf unterbrochen.}

\subsection{Planungsunterstützung, Szenarien und Werkzeugökosystem}

\textbf{Interviewer:} When you describe the VCDM as a planning-support tool, who is it intended for? Is it aimed at particular users, such as municipal planners, or at any stakeholder involved in making decisions during the process?

\textbf{Raoul Bunschoten:} It can be used by any stakeholder. For example, I am part of a group called Zukunft Berlin. We recently celebrated its twentieth anniversary with the mayor and state secretaries. That setting includes stakeholders at the political level: the people governing Berlin and instructing state-owned housing institutions regarding what kind of housing to build and which materials to use. Politicians are therefore also stakeholders.

I am meeting representatives of the Senate and Tegel Projekt concerning a phase of the Schumacher Quartier. They have a mandate to build in wood, but implementing that mandate is proving very complicated. We are working on how to move the process forward. The relevant stakeholders can therefore change from one situation to another.

At other times, it is possible to assemble representatives from across the complete value chain. A few years ago, we organised a series of funded workshops that achieved this. We presented the VCDM, although the participants did not directly operate the digital model. Instead, we used analogue games and gamification methods, similar to the techniques used at the chair, to help people begin communicating with one another.

In this sense, the VCDM can be understood as an interface that supports analogue processes of navigation and collaborative planning. That is why I call it a planning-support tool.

Many people imagine that the tool will solve a problem for them or generate substantial financial returns. We have to say: no, it does not make the decisions. Stakeholders still have to decide for themselves. The tool creates values and evidence that can be compared. In German, it is a \textit{Bewertungsinstrument}—an evaluation instrument.

You can create scenarios and then support those scenarios with evidence. We use open-source data and hundreds of data sources. Satellite data, for example, allows us to identify trees across Brandenburg, understand where wood is located, and determine who owns it. The VCDM is therefore an interface with access to datasets, although what we currently have is still a prototype.

The principle is that different facts can be placed alongside one another and fed into a simulation. The VCDM is a family of tools, not one tool. Other tools connect to it. The VCDM acts as a platform, while additional tools are developed for particular parts of the value chain. Further development depends partly on funding. The resulting system is more complex than it might initially appear and becomes a form of human-machine interactive interface.

An important issue in discussions about digital twins is the need to capture dynamics. The goal is not a still image but a dynamic image. This relates again to \textit{Urban Curation}, which is concerned with managing dynamic planning situations.

\subsection{Dynamische Planung und unvollständige Daten}

\textbf{Interviewer:} What would be a practical example of a “dynamic image”? Are you referring to real-time data changes, or to something else?

\textbf{Raoul Bunschoten:} Consider a forest. When you see a point cloud, you may assume it is a complete representation of what is there. However, a point cloud does not necessarily capture everything. It may fail to represent trees that have already been cut and are lying on the ground. Those trees might be recorded as dead even though the material remains in good condition. These hybrid situations require interpretation. Pure data do not tell the full story.

Experts are therefore needed throughout the value chain to interpret data and compare their observations. They must be able to return to the information and say, “Let us compare notes.”

A different example is the changing situation in housing and construction regulation. The \textit{Bau-Turbo} is intended to accelerate building by simplifying certain regulatory processes. The \textit{Holzbau-Richtlinie} provides guidance for timber construction but is not itself a law. Developments of this kind affect projects that are already under way. This is what I mean by dynamics: the project has to adapt.

Suppose an architect has completed a design and is being paid for that work, but the regulations suddenly change. We are currently dealing with a situation of this kind in the Schumacher Quartier, although I cannot discuss the details. A project may have been designed and planned, and then a regulatory or planning condition changes.

For example, the required dimensions of a laminated beam may be affected by its span. The design may appear precise and resolved, but a new regulatory condition can make the height of the beam problematic. The beam may still be structurally capable of doing the work, but another requirement changes the situation. That is the type of dynamic condition the model must help stakeholders understand.

\textbf{Interviewer:} So the VCDM is not a design tool in the conventional sense. It is more like a collaborative interface that guides and supports planning.

\textbf{Raoul Bunschoten:} That is correct. It helps you plan, but it also helps guide you through the stages of an implementation process.

\subsection{Baukastensysteme und Skalierung}

\textbf{Interviewer:} How does the platform relate to the idea of building kits? At what stage of design do these kits become relevant, who designs them, and how do they connect to the platform?

\textbf{Raoul Bunschoten:} We work with architects as well as industrial companies that manufacture building components. The idea of a building kit emerges from asking where meaningful impact can be achieved.

You can work from the bottom up by developing technologies, creating an intelligent design, applying BIM, or using AI to test combinations. That is one direction. But you can also work from the top down. If the government says that 400,000 housing units are required, you must ask industry whether it can deliver that volume. No single industry can produce the complete volume required by such a strategy.

You therefore turn to organisations such as Fraunhofer and ask how multiple industries might be combined through distributed production. That could enable a transition from business as usual toward bio-technic construction.

This includes both \textit{Bauen im Bestand}—working with existing buildings—and \textit{Neubau}, or new construction. We do not treat them as mutually exclusive. There is a strong movement toward working with existing buildings, and your research is connected to that. But the situation is more complicated than simply declaring that no new building should occur. Many people in Berlin need housing, and they need it quickly. It cannot be framed as an either-or choice.

The digital model itself does not possess a political or moral preference. However, the supply chains and datasets for existing buildings and new construction are different. Once you consider goals such as 400,000 units nationally or 220,000 units in Berlin, you have to think about production differently.

The question becomes whether a form of mass production can be created while avoiding the limitations traditionally associated with standardisation. We know mass production from systems such as the \textit{Plattenbau}, which has been both criticised and successful in different contexts. Advanced industries already mass-produce products such as cars and phones.

Can construction achieve large numbers while applying the agile and ambidextrous methods of Industry 4.0? This does not mean that every unit must be identical. Standard components can be defined digitally and organised into component libraries. Those libraries can include newly manufactured components as well as elements that already exist and can be reused.

Certain typologies may have to be agreed upon, but most housing typologies are not completely different from one another. Many architects can work with a shared set of basic elements while still composing different solutions.

Automation also enables variation. A robot can produce ten housing units with small differences without the difficulty that such variation might create in a conventional repetitive process. With Fraunhofer, we are also examining cobots and the automation of construction. This is where the idea of the building kit becomes useful.

Building kits have a long history. The German term is \textit{Baukasten}. Their basic purpose in this context is to enable upscaling while retaining adaptability.

Upscaling also connects to broader debates about urban growth. Some people argue against further urban development and suggest that the city is already built. I do not share such a simple opposition. Society is dynamic: needs change, people move, and cities adapt. Some form of growth will continue to occur. Urban growth may also contribute to carbon storage if bio-based materials are used effectively.

I therefore do not regard the growth of cities as inherently negative. I am also involved in work concerning the regeneration of forestry in Indonesia and Thailand. Forestry and urban construction should be understood as parts of the same value chain, not as opposites.

\subsection{Regulierung, Szenariotechniken und synthetische Daten}

\textbf{Interviewer:} How do building regulations enter the digital platform, and how are they connected to the building kits?

\textbf{Raoul Bunschoten:} That is a good and difficult question. At a technical level, the information has to be tagged and introduced into the tool. I am not the technical specialist, so members of the research team would explain the exact method. At a more general level, however, regulatory information has to be entered and interpreted.

We sometimes translate this information into very simple controls. The VCDM is an interactive simulation tool that can use sliders, for example. Certain sliders can be linked to the numerical values contained in a regulation. At some point, however, the slider reaches a limit. When the numerical model can go no further, the process returns to analogue negotiation.

The scenario techniques developed over the years through CHORA and at the university are intended to support this transition between model and negotiation. Years ago, we participated in a major EU call through a project called Citizen City Maker. We worked with two AI companies on the relationship between scenarios created by people and simulations that used numerical values derived from those scenarios.

That work was extremely difficult at the time because language models were not yet capable of extracting and processing the required information from scenarios. They are more capable now. Another difficulty was the limited quantity of data that could be derived from each scenario.

We therefore looked at methods used in fields such as medical and cancer research, where artificial or synthetic data are created. We began producing artificial data as a parallel research process several years ago. I find this very exciting, although it goes somewhat beyond the present discussion.

The question is how large language models can generate and work with highly dimensional data. That is connected to the research you are doing. I would be very interested to see what you are producing for Zukunft Bau. It took a long time for your project to receive support, and the funding itself is not especially large, but I consider the project very strong. I still have the project documents and the book. We should take time to examine the work together.

\subsection{Vom akademischen Prototyp zur Planungsinfrastruktur}

\textbf{Interviewer:} My final question is: what would be required for a platform such as the VCDM to move from an academic prototype into everyday planning infrastructure?

\textbf{Raoul Bunschoten:} That is a very good question, and the transition is extremely difficult. At present, we have prototypes. VCDM 2 is highly simplified and remains a raw prototype. Systems of this scale can only remain broad prototypes unless several million euros are available for development. The technology economy in which these systems operate deals in billions.

Nevertheless, we are in contact with several Berlin districts and with organisations such as HOWOGE, the state-owned housing companies, and other institutions responsible for tens of thousands of housing units.

We are particularly active in the issue known in Germany as \textit{Aufstockung}: using existing housing to create opportunities for vertical extensions. Building kits can also be applied to gaps in the urban fabric. Students have developed intelligent projects examining how such gaps can be scanned and addressed at scale.

At the same time, housing associations still have to produce a certain number of units on new sites. New construction remains a real demand, even when it is politically contested. Our claim is that adaptable building kits can be developed for vertical extensions, urban infill, and new sites simultaneously. None of these applications is purely industrial or completely identical.

People in the Netherlands and elsewhere in the European Union are working on related ideas, and we are in contact with several of them. A tool of this kind can be a strong negotiation instrument, but the interface must also be simplified.

We are currently working on two small commissions connected to something developed at the beginning of Bio-Tech 4.0: the \textit{Holzbau-Matrix}. We designed an initial evaluation matrix for Tegel Projekt, connected to Tegel TXL. We are now beginning to redevelop it for a hybrid application.

This is a much simpler form of interface. Tools used in procurement do not necessarily need to resemble dashboards or contain visible sliders. They may become very simple systems showing what exists where and how different options are scored. At that point, relatively little visible technology remains. The underlying thinking still comes from the VCDM, but the practical interface is reduced to what is useful in a procurement process.

We are also discussing workshops with districts and other organisations. A few months ago, we held a workshop in which members worked through two VCDM scenarios. It was interesting to observe how the process functioned. Financial actors were also involved, including representatives connected to Madaster. This raises the question of how a banker or financing institution engages with the VCDM.

We are gradually exploring how these actors can be involved in procurement and planning processes. It is a difficult and slow process. Some people resist the technology and ask why it is necessary. Others misunderstand the system as a machine that automatically generates solutions or as a simple digital twin or mapping tool. The terminology and definitions are still unsettled: people are not always clear about what they are dealing with.

I am therefore increasingly concerned with the professional role required to work across the overall value chain. Throughout my work at the chair, I have asked where the planners are and how planners can be trained to operate creatively. How can urban planning be treated as an artistic practice and as part of the creative industries?

The conceptual question is: who actually performs this work? Who has the expertise to speak with an architect, a housing association, and a banker, and to negotiate through regulations and financial decisions until a project can move forward? Many processes are currently blocked because that coordinating expertise and professional role are not sufficiently established.

Berlin is under pressure to build housing very quickly. In that pressure, bio-technic construction and climate objectives can disappear from the discussion, which risks creating another crisis. I am interested in how we can define a professional role that is socially recognised, can be formally commissioned, and can eventually be certified. Developing that role is one of the projects I have set myself in my emeritus work.

Our aim is also to develop a third stage of the VCDM that includes the issues I have described. We are waiting for additional sponsorship and funding, although some resources have already been allocated.

I would encourage you to continue your work because the challenge is enormous. People sometimes say that other organisations, such as Madaster, are already doing the same thing. We have to be careful about treating this field primarily as a competition. The task is too large. At some point, the different initiatives need to meet, cooperate, and connect their work. Keep going, and let us continue the conversation.

